\guidelinesubsubsection{Rationale}

Reproducibility depends on access to the model under study.
When research relies exclusively on proprietary models, other researchers cannot independently verify or build upon the findings.
Including an open LLM as a baseline ensures that at least part of the study can be fully replicated.

\guidelinesubsubsection{Recommendations}

Empirical studies using LLMs in SE, especially those that target commercial tools or models, \should incorporate an open LLM as a baseline and report established metrics for inter-model agreement (see \benchmarksmetrics).
We acknowledge that including an open LLM baseline might not always be possible, for example, if the study involves human participants, and letting them work on the tasks using two different models might not be feasible.
Using an open model as a baseline is also not necessary if the use of the LLM is tangential to the study goal.

Open models allow other researchers to verify research results and build upon them, even without access to commercial models.
A comparison of commercial and open models also allows researchers to contextualize model performance.
Researchers \should ensure the open-LLM baseline is independently reproducible from their \supplementarymaterial.

Open LLMs are available from hubs such as \href{https://huggingface.co/}{\emph{Hugging Face}}. They can be self-hosted with frameworks such as \href{https://ollama.com/}{\emph{Ollama}} or \href{https://lmstudio.ai/}{\emph{LM Studio}}, accessed through cloud services such as \href{https://together.ai/}{\emph{Together AI}}, AWS, Azure, and Google Cloud, or routed through aggregators such as \href{https://openrouter.ai/}{\emph{OpenRouter}} that expose many providers behind a single API. For agentic setups, open-source tools such as \href{https://www.continue.dev/}{\emph{Continue}}, \href{https://cline.bot/}{\emph{Cline}}, and \href{https://opencode.ai/}{\emph{opencode}}~\cite{opencode} publish their full agent code, system prompts, and tool catalog. By contrast, vendor-hosted services such as \emph{GitHub Copilot} and \emph{Claude Code} expose only parts of their tooling (e.g., editor extensions, hook examples) and keep their deployed agent loops, system prompts, and tool catalogs proprietary.
% GROUNDING opencode: "OpenCode is an open source AI coding agent."

The term ``open'' can have different meanings in the context of LLMs.
\citeauthor{widder2024open} discuss three types of openness (i.e., transparency, reusability, and extensibility) and what openness can and cannot provide~\cite{widder2024open}.
% GROUNDING widder2024open: "We highlight three main affordances of 'open' AI, namely transparency, reusability, and extensibility"
The \emph{Open Source Initiative} (OSI)~\cite{OSIAI2024} defines open-source AI as having access to everything needed to understand, modify, share, retrain, and recreate the model.
% GROUNDING OSIAI2024: "Sufficiently detailed information about the data used to train the system so that a skilled person can build a substantially equivalent system."

\guidelinesubsubsection{Examples}

An increasing number of studies have adopted open LLMs as baseline models.
For example, \citeauthor{DBLP:journals/corr/abs-2406-09834} evaluated seven advanced LLMs, six of which were open-source, testing 145 API mappings drawn from eight popular Python libraries across 28,125 completion prompts aimed at detecting deprecated API usage in code completion~\cite{DBLP:journals/corr/abs-2406-09834}.
% GROUNDING DBLP:journals/corr/abs-2406-09834: "This study involved seven advanced LLMs, 145 API mappings from eight popular Python libraries, and 28,125 completion prompts."
\citeauthor{DBLP:journals/corr/abs-2408-04430} compared four LLMs on a cross-language code clone detection task~\cite{DBLP:journals/corr/abs-2408-04430}.
% GROUNDING DBLP:journals/corr/abs-2408-04430: "We evaluated four LLMs and one embedding model for cross-lingual code clone detection."
Three evaluated models were open-source.
\citeauthor{DBLP:conf/eicc/GoncalvesSC0MPS25} fine-tuned the open LLM \emph{LLaMA} 3.2 on a refined version of the \emph{DiverseVul} dataset to benchmark vulnerability detection performance~\cite{DBLP:conf/eicc/GoncalvesSC0MPS25}.
% GROUNDING DBLP:conf/eicc/GoncalvesSC0MPS25: "we present a refined version of DiverseVul dataset, which is used to fine-tune a large language model, LLaMA 3.2, for vulnerability detection"
\citeauthor{DBLP:journals/corr/abs-2601-11895} evaluated nine code completion models on the \emph{DevBench} benchmark and included three open-weight models (\emph{DeepSeek-V3}, \emph{DeepSeek-V3.1}, and \emph{Ministral-3B}) alongside commercial frontier models, releasing benchmark, evaluation scripts, and per-model raw completions under an MIT license~\cite{DBLP:journals/corr/abs-2601-11895}.
% GROUNDING DBLP:journals/corr/abs-2601-11895: "Claude 4 Sonnet, Claude 3.7 Sonnet, GPT-4.1 mini, GPT-4.1, GPT-4o, DeepSeek-V3, DeepSeek-V3.1, GPT-4.1 nano, Ministral 3B"
\emph{CodeBERT}, a bimodal transformer pre-trained by Microsoft Research, is published with model weights, source code, and data processing scripts on GitHub~\cite{codebert}. It has been used as an open baseline across diverse SE tasks, including exploit code generation~\cite{DBLP:journals/jss/YangZCZHC23}, vulnerability detection~\cite{DBLP:conf/gaiis/XiaSD24}, code clone detection~\cite{DBLP:conf/kbse/SonnekalbGBM22}, and programming assistance for exception handling~\cite{DBLP:conf/icse/CaiYMMN24}.
% GROUNDING codebert: "This repo provides the code for reproducing the experiments in CodeBERT: A Pre-Trained Model for Programming and Natural Languages."
% GROUNDING DBLP:journals/jss/YangZCZHC23: "we propose a novel template-augmented exploit code generation approach ExploitGen based on CodeBERT"
% GROUNDING DBLP:conf/gaiis/XiaSD24: "this study introduces a software defect detection system leveraging CodeBERT and Bi-LSTM"
% GROUNDING DBLP:conf/kbse/SonnekalbGBM22: "Transformer networks such as CodeBERT already achieve outstanding results for code clone detection in benchmark datasets"
% GROUNDING DBLP:conf/icse/CaiYMMN24: "we design Neurex via a multi-tasking model with the fine-tuning of the large language model CodeBERT for the three above exception handling recommending tasks"

\guidelinesubsubsection{Benefits}

A true open LLM baseline improves reproducibility by exposing model architectures, parameter settings, and ideally training data, enabling independent verification of results.
Such baselines also let researchers compare novel methods against a stable reference point, since proprietary models can silently change between tests.
They also allow inspection of training data (when released) and model behavior, helping identify biases and limitations.
Unlike closed-source alternatives, which can be withdrawn or silently updated, open LLMs remain available for future studies.
They typically avoid the per-use API fees that can constrain research groups with limited budgets.

\guidelinesubsubsection{Challenges}

Open-source LLMs face several challenges:
\begin{itemize}
    \item \emph{Definitional inconsistency.} Many models release only trained weights without training data or methodological details (``open weight'' openness)~\cite{Gibney2024}, which is why we reference the OSI definition in our recommendations~\cite{OSIAI2024}.
    % GROUNDING Gibney2024: "Technology giants such as Meta and Microsoft are describing their artificial intelligence (AI) models as 'open source' while failing to disclose important information about the underlying technology"
    % GROUNDING OSIAI2024: "The complete source code used to train and run the system"
    \item \emph{Performance gap.} Open models often lag behind the most advanced proprietary ``frontier'' models in common benchmarks, making it difficult to demonstrate clear improvements when evaluating new methods using open LLMs alone.
    \item \emph{Hardware demands.} Deploying and experimenting with these models typically requires substantial hardware resources, in particular high-performance GPUs that may be beyond reach for many academic groups.
    \item \emph{Operational complexity.} Unlike APIs provided by proprietary vendors (e.g., the OpenAI API), installing, configuring, and fine-tuning open-source models can be technically demanding.
\end{itemize}

\guidelinesubsubsection{Study Types}

This guideline applies primarily to study types in which the researcher controls which LLM is used.
For \benchmarkingtasks, an open LLM \should be one of the models under evaluation, so that the reported scores can be independently re-run. Where this is not feasible, researchers \should justify the omission and, per \design, ensure the evaluation harness can be used with open models.
For controlled experiments under \llmusage, an open LLM \should be one of the models under test; if the experimental design requires capabilities that only a specific commercial model exhibits, researchers \should acknowledge this as a limitation.
When evaluating \newtools, researchers \should use an open LLM as a baseline whenever it is technically feasible; if integration proves too complex, they \should report the initial benchmarking results of open models.
For \annotators and \judges, researchers \should compare annotation or judgment quality from open vs.\ commercial models to assess the extent to which results depend on a specific proprietary model.
For \synthesis, researchers \should compare synthesis results from open and commercial models to evaluate robustness of the findings.
For \llmusage, using an open LLM as a baseline is often not feasible when the study observes participants using specific commercial tools; in such cases, investigators \should explicitly acknowledge its absence and discuss how this limitation might affect their conclusions.
For \subjects, using an open LLM as a baseline may similarly be impractical if the study design requires the capabilities of a specific model; researchers \should acknowledge this limitation when applicable.

\guidelinesubsubsection{Advice for Reviewers}

Reviewers should distinguish between LLM use that is central to the research (e.g., building LLM-driven SE tools, using LLMs for synthesis) and use that is tangential (e.g., generating recruiting materials). An open LLM baseline is expected only when LLM use is central; otherwise, its absence need not be justified.
When an open baseline is expected, reviewers should look for either the use of an open LLM or a convincing argument for why it is impractical. If authors claim openness, some justification for that characterization is appropriate. Where practical, reviewers should examine whether the replication package contains sufficiently detailed instructions.
Performance differences between open and proprietary models are beyond the study authors' control. Reviewers should focus on whether the open baseline serves the study's methodological purpose rather than on its absolute performance.

\guidelinesubsubsection{See Also}
\begin{itemize}[label=$\rightarrow$]
    \item \seemodelversion: Authors must report version, configuration, and parameters for open models, not just commercial ones.
    \item \seedesign: Self-hosting an open model adds hardware, hosting, and infrastructure to document.
    \item \seebenchmarksmetrics: Inter-model agreement metrics make open-vs-commercial comparisons interpretable.
    \item \seelimitationsmitigations: When using an open LLM as a baseline is impractical, authors must report this as a limitation.
\end{itemize}
